\section{Große Anwartschaftsversicherung}\label{sec:awv}

Für Versicherte mit ruhendem Leistungsanspruch (z.\,B. vorübergehende
Versicherungspflicht in der GKV) kann der Vertrag als große
Anwartschaftsversicherung fortgeführt werden. Der Anwartschaftsbeitrag
deckt den Barwert des Alterungsrückstellungsaufbaus sowie das Recht auf
Wiederaufleben des Versicherungsschutzes ohne erneute Risikoprüfung.
Die formelmäßige Herleitung der Anwartschaftsbeiträge aus den
Grundtarifbeiträgen (Prozentsätze der Beitragsraten des
Abschnitts~\ref{sec:beitragsraten}), die Zuschläge sowie die
Bilanzdeckungsrückstellung der Anwartschaftsversicherung sind in der
Rahmen-TB \textbf{RTB-AWV (Version 2.1)} dokumentiert
(Abschnitt~\ref{sec:rahmentb}); die dortigen Bestimmungen gelten für
den Tarif VitalPlus unverändert. Für die Zuschläge gilt
Abschnitt~\ref{sec:zuschlaege} entsprechend.

\section{Kleine Anwartschaftsversicherung}\label{sec:kleineawv}

Für die kleine Anwartschaftsversicherung (Erhalt der
Alterungsrückstellung ohne Aufbau weiterer Rückstellungen) gilt
Abschnitt~\ref{sec:awv} entsprechend; Jahresnettobeiträge, Zuschläge
und Beitragsraten sind ebenfalls in der Rahmen-TB \textbf{RTB-AWV
(Version 2.1)} dokumentiert.
\clearpage
